%-------------------------------------------------------------------------------
\section{Non-cohesive suspended sediment transport}
\label{sec:NCO:susp}
%-------------------------------------------------------------------------------
%2D suspended transport
In 2D cases, the suspended sediment transport is accounted by solving the two-dimensional advection-diffusion equation, expressed by:
\begin{equation}\label{eq:2DADE}
\frac{\partial hC}{\partial t} + \frac{\partial hUC}{\partial x} + \frac{\partial hVC}{\partial y} =
\frac{\partial}{\partial x}\left(h\epsilon_s\frac{\partial C}{\partial x}\right) +
\frac{\partial}{\partial y}\left(h\epsilon_s\frac{\partial C}{\partial y}\right) + E-D
\end{equation}
where $C=C(x,y,t)$ is the depth-averaged concentration \textcolor{black}{\bf expressed in g/l}, $(U,V)$ are the depth-averaged components of the velocity in the $x$ and $y$ directions, respectively, $\epsilon_s$ is the turbulent diffusivity of the sediment, often related to the eddy viscosity $\epsilon_s=\nu_t/\sigma_c$, with $\sigma_c$ the Schmidt number. In our case, $\sigma_c=1.0$.

%3D suspended transport
In 3D cases, the suspended sediment transport is accounted by solving the three-dimensional advection-diffusion equation, expressed by:
\begin{equation}\label{eq:3DADE}
\frac{\partial C}{\partial t} + \frac{\partial uC}{\partial x} + \frac{\partial vC}{\partial y} + \frac{\partial wC}{\partial z} -\frac{\partial w_sC}{\partial z} =
\frac{\partial}{\partial x}\left(\epsilon_s\frac{\partial C}{\partial x}\right) +
\frac{\partial}{\partial y}\left(\epsilon_s\frac{\partial C}{\partial y}\right) +
\frac{\partial}{\partial z}\left(\epsilon_s\frac{\partial C}{\partial z}\right)
\end{equation}
where $C=C(x,y,z,t)$ is the concentration \textcolor{black}{\bf expressed in g/l}, $(u,v,w)$ are the components of the velocity in the $x$, $y$ and $z$ directions, respectively. $w_s$ is the sediment settling velocity. $\epsilon_s$ is the turbulent diffusivity of the sediment, often related to the eddy viscosity $\epsilon_s=\nu_t/\sigma_c$, with $\sigma_c$ the Schmidt number. In our case $\sigma_c=1.0$.

%...............................................................................
\subsection{Initial conditions for suspended sediment transport}
%...............................................................................
The initial condition value of the concentration can be specified through the keyword \telkey{INITIAL SUSPENDED SEDIMENTS CONCENTRATION VALUES} (real type list, {\ttfamily = 0.0} by default), following the order given in \telkey{CLASSES TYPE OF SEDIMENT}.
These values must be \textcolor{black}{\bf expressed in g/l}. In 2D, this keyword is not considered on boundary nodes if \texttt{EQUILIBRIUM INFLOW CONCENTRATION = YES}.\\

\begin{lstlisting}[frame=trBL]
INITIAL SUSPENDED SEDIMENTS CONCENTRATION VALUES :0.D0;1.D0;...
\end{lstlisting}

\subsection{Boundary conditions for suspended sediment transport}

The specification of boundary conditions is done in a boundary condition file, usually named with extension \texttt{*.cli}. The reader is referred to \S\ref{sec:flags} for the definition of the different flags used in the boundary condition file.

\subsubsection{Wall boundary conditions}
At banks and islands, the no-flux boundary condition is imposed: the suspended load concentration gradients are set to zero.
\begin{equation}
\label{eq:BC-noflux}
\epsilon_s\frac{\partial C}{\partial n}=\epsilon_s \Grad C \cdot \vec{n}=0
\end{equation}
For this case, the flag \texttt{LIEBOR} is set \texttt{= 2} as shown in the example below:

\begin{lstlisting}[frame=trBL]
2 2 2 0.0 0.0 0.0 0.0 (*@\color{PantoneRed}2@*) 0.0 0.0 0.0 565 1
\end{lstlisting}

\subsubsection{Inflow boundary conditions}
Several ways to impose the concentration at inflow boundary conditions are proposed in \gaia. For this case, the flag \texttt{LICBOR} is set \texttt{= 5}, with the following options for the concentration values:
\begin{itemize}
\item Specified in the column \texttt{CBOR}. In the example below, a concentration value {\ttfamily = 1.0} g/l is provided:
\begin{lstlisting}[frame=trBL]
4 5 5 0.0 0.0 0.0 0.0 (*@\color{PantoneRed}5@*) (*@\color{PantoneRed}1.0@*) 0.0 0.0 565 1
\end{lstlisting}
\item Specified by the concentration values at the boundary through the keyword \telkey{PRESCRIBED SUSPENDED SEDIMENTS CONCENTRATION VALUES} (real type list of size equal to the number of boundaries $\times$ the number of sediment classes), expressed in g/l:
\begin{lstlisting}[frame=trBL]
4 5 5 0.0 0.0 0.0 0.0 (*@\color{PantoneRed}5@*) 0.0 0.0 0.0 565 1
\end{lstlisting}
\begin{lstlisting}[frame=trBL]
  PRESCRIBED SUSPENDED SEDIMENTS CONCENTRATION VALUES =
  0.0; 0.0; 1.0; ...
\end{lstlisting}
The order is the following: boundary 1 (class 1, class2, etc.), then boundary 2, etc.
\item Computed by \gaia{} in 2D cases through the keyword \telkey{EQUILIBRIUM INFLOW CONCENTRATION} (logical type variable, {\ttfamily = NO} by default) and according to the choice of the formula of equilibrium near-bed concentration (\telkey{SUSPENSION TRANSPORT FORMULA FOR ALL SANDS = 1} by default)

\begin{lstlisting}[frame=trBL]
4 5 5 0.0 0.0 0.0 0.0 (*@\color{PantoneRed}5@*) 0.0 0.0 0.0 565 1
\end{lstlisting}
\begin{lstlisting}[frame=trBL]
EQUILIBRIUM INFLOW CONCENTRATION = YES
SUSPENSION TRANSPORT FORMULA FOR ALL SANDS = 1
\end{lstlisting}
\item Time-varying concentration values (in g/l) must be provided in an external \texttt{ASCII} file. In this case the keyword \telkey{LIQUID BOUNDARIES FILE} has to be used in the hydrodynamic steering file of \telemac{2D} or \telemac{3D} according to the case. Suspended sediment concentrations are treated like tracers, the reader is thus referred to the user manual of modules \telemac{2D} or \telemac{3D}.
If tracers and sediments are declared at the same time, sediments will be allocated after tracers. For example if we have 3 tracers and 2 suspended sediments, the first sediments will be the fourth tracer.
\begin{lstlisting}[frame=trBL]
4 5 5 0.0 0.0 0.0 0.0 (*@\color{PantoneRed}5@*) 0.0 0.0 0.0 565 1
\end{lstlisting}
\begin{lstlisting}[frame=trBL]
LIQUID BOUNDARIES FILE = '<file name>'
\end{lstlisting}
\item Vertical profile provided in 3D cases through the keyword \telkey{VERTICAL PROFILES OF SUSPENDED SEDIMENTS}, which can be equal to:
\begin{itemize}
\item 0 if the profile is programmed by users,
\item 1 if constant profile,
\item 2 if Rouse profile,
\item 3 if normalised Rouse profile and assigned concentration,
\item 4 if modified Rouse profile with molecular viscosity.
\end{itemize}
\end{itemize}

According to the hydrodynamic forcing, the use of the keyword \telkey{EQUILIBRIUM INFLOW CONCENTRATION} can be used to prevent the excessive erosion and deposition on the inflow boundary (for 2D cases).

\subsection{Outflow boundary conditions}
At the outflow boundary, the suspended load concentration gradient in the flow direction is set to zero. For this case, the flag \texttt{LICBOR} is set \texttt{= 4} as shown in the example below:

\begin{lstlisting}[frame=trBL]
5 4 4 0.0 0.0 0.0 0.0 (*@\color{PantoneRed}4@*) 0.0 0.0 0.0 565 1
\end{lstlisting}

Stability issues at inflow boundaries can be solved by activating the keyword \telkey{TREATMENT OF FLUXES AT THE BOUNDARIES} (integer type variable, {\ttfamily = 2}, default option) in the steering file of \telemac{2D} or \telemac{3D}. The choice {\ttfamily = 1} can be activated when prescribed values are provided. The choice {\ttfamily = 2} is used when fluxes are imposed.

%...............................................................................
%\subsection{Diffusion and dispersion} -> what about this part now?
%%...............................................................................
%The diffusion term in the ADE for the depth-averaged suspended concentration is taken into account with the keyword \telkey{ DIFFUSION} (logical type, set {\ttfamily = YES} by default). The values of the diffusion coefficients can be specified with the keyword \telkey{ OPTION FOR THE DISPERSION} (integer type, set {\ttfamily = 1} by default), with the following options:
%\begin{itemize}
%\item {\ttfamily = 1}: values of the longitudinal and transversal dispersion coefficients are provided with the keywords \texttt{DISPERSION ALONG THE FLOW} and \texttt{DISPERSION ACROSS THE FLOW}, respectively:
%\begin{itemize}
%\item \telkey{DISPERSION ALONG THE FLOW} (real type, set {\ttfamily = $1.0 \times 10^{-2}$ m$^2/$s} by default)
%\item \telkey{DISPERSION ACROSS THE FLOW} (real type, set {\ttfamily = $1.0 \times 10^{-2}$ m$^2/$s} by default)
%\end{itemize}
%
%\item {\ttfamily = 2}: values of the longitudinal and transversal dispersion coefficients are computed with the Elder model $T_l=\alpha_l u_* h$ and $T_t=\alpha_t u_* h$, where the coefficients $\alpha_l$ and $\alpha_t$ can be provided with the keywords \texttt{DISPERSION ALONG THE FLOW} and \texttt{DISPERSION ACROSS THE FLOW}.
%\item {\ttfamily = 3}: values of the dispersion coefficients are provided by the hydrodynamics module (e.g. \telemac{2D})
%\end{itemize}

%...............................................................................
\subsection{Numerical treatment of the diffusion terms}
%...............................................................................
The keyword \texttt{OPTION FOR THE DIFFUSION OF TRACER} (integer type, set {\ttfamily = 1} by default) allows to choose the treatment of the diffusion terms in the advection-diffusion equation \ref{eq:2DADE} for the depth-averaged suspended concentration:
\begin{itemize}
\item \texttt{= 1}: the diffusion term is solved in the form $\nabla\cdot(\varepsilon_s\nabla T)$
\item \texttt{= 2}: the diffusion term is solved in the form $\frac{1}{h}\nabla\cdot(h\varepsilon_s\nabla T)$
\end{itemize}
This keyword must be activated in the steering file of \telemac{2D}; user can refer to the corresponding user manual for further details.
%...............................................................................
\subsection{Numerical treatment of the advection terms}
%...............................................................................
The choice for the scheme for the treatment of the advection terms can be done with the keyword \texttt{SCHEME FOR ADVECTION OF SUSPENDED SEDIMENTS} (integer type, set {\ttfamily = 5} by default):
\begin{lstlisting}[frame=trBL]
1="CHARACTERISTICS"
2="SUPG"
3="CONSERVATIVE N-SCHEME"
4="CONSERVATIVE N-SCHEME"
5="CONSERVATIVE PSI-SCHEME"
13="EDGE-BASED N-SCHEME"
14="EDGE-BASED N-SCHEME"
15="ERIA SCHEME"
\end{lstlisting}

The keyword \texttt{SCHEME OPTION FOR ADVECTION OF SUSPENDED SEDIMENTS} can be used to choose futher options for advection schemes. In particular, if characteristics are used, the following options are possible :
\begin{itemize}
\item \texttt{1}: strong form;
\item \texttt{2}: weak form.
\end{itemize}
If advection is solved by the N or PSI scheme, options are:
\begin{itemize}
\item \texttt{1}: explicit scheme;
\item \texttt{2}: first order predictor-corrector scheme;
\item \texttt{3}: second order predictor-corrector scheme;
\item \texttt{4}: local semi-implicit scheme.
\end{itemize}
To have more details about these keywords, readers can refer to the \telemac{2D} or \telemac{3D} user guide (cf. chapters about tracer advection). In fact, the advection scheme used by sediments is exactly the same used by tracers, implemented in hydrodynamics modules. Keywords \texttt{SCHEME FOR ADVECTION OF SUSPENDED SEDIMENTS} and \texttt{SCHEME OPTION FOR ADVECTION OF SUSPENDED SEDIMENTS} have been added from V8P2 in \gaia{} in order to have proper steering files.

\begin{WarningBlock}{Note:}
  It is recommended to use the schemes {\ttfamily 4} or {\ttfamily 14} for a good compromise between accuracy and computational time (specially if tidal flats are present). It is also suggested to activate
  the keyword \texttt{CONTINUITY CORRECTION = YES}.
\end{WarningBlock}

A brief description of the numerical schemes implemented in \gaia{} is given below:
\begin{itemize}
\item \textbf{Method of characteristics} (\texttt{1})
\begin{itemize}
\item Unconditionally stable and monotonous
\item Diffusive for small time steps, Not conservative
\end{itemize}
\item \textbf{Method Streamline Upwind Petrov Galerkin SUPG} (\texttt{2})
\begin{itemize}
\item Based on the Courant number criteria
\item Less diffusive for small time steps, Not conservative
\end{itemize}
\item \textbf{Conservative N-scheme (similar to finite volumes)} (\texttt{3, 4})
\begin{itemize}
\item Solves the continuity equation under its conservative form
\item Recommended for correction on convection velocity
\item Courant number limitation (sub-iterations to reduce time step)
\end{itemize}
\item \textbf{Edge-based N-scheme NERD} (\texttt{13, 14})
\begin{itemize}
\item Same as \texttt{3} and \texttt{4} but adapted to tidal flats
\item Based on positive-depth algorithm
\end{itemize}
\item \textbf{Distributive schemes PSI} (\texttt{5})
\begin{itemize}
\item fluxes corrected according to the tracer value: relaxation of Courant number criteria, less diffusive than
schemes \texttt{4, 14} but larger CPU time
\item Should not be applied for tidal flats
\end{itemize}
\item \textbf{Eria scheme} (\texttt{15})
\begin{itemize}
\item Works for tidal flats
\end{itemize}
\end{itemize}
For further information about these schemes, refer to \telemac{2D} user manual.

\subsection{Correction of the convection velocity}
As most of the sediment transport processes occur near the bed, it often exhibits an over-estimation of suspended sand transport.
A correction method accounting for the vertical velocity and concentration profiles is therefore proposed. A straightforward treatment of the advection terms would imply the
definition of an advection velocity and replacement of the depth-averaged
velocity $U$ along the $x-$axis in Eq.~(\ref{eq:2DADE}) by:
\begin{equation*}
U_{conv} = \overline{UC}/C.
\end{equation*}
A correction factor is introduced in \gaia, defined by:
\begin{equation*}
F_{conv} =\frac{U_{conv}}{U}.
\end{equation*}
A similar treatment is done for the depth-averaged velocity $V$ along the $y-$axis. For further details, see~\cite{Huybrechts}.
The convection velocity should be smaller than the mean flow velocity ($F_{conv} \leq 1$) since sediment concentrations are mostly transported in the lower part of the water column where velocities are smaller. We further
assume an exponential concentration profile which is a reasonable
approximation of the Rouse profile, and a logarithmic velocity profile, in
order to establish the following analytical expression for $F_{conv}$:
\begin{equation*}
F_{conv} =-\frac{I_2 - \ln\left(\frac{B}{30}\right) I_1}{I_1 \ln\left(
\frac{eB}{30}\right)},
\end{equation*}
with $B=k_s/h = Z_{ref}/h$ and
\begin{equation*}
I_1 =\int_B^1\left(\frac{(1-u)}{u}\right)^R du,\quad I_2 = \int_B^1 \ln u \left(\frac{(1-u)}{u} \right)^R du.
\end{equation*}

The keyword \texttt{CORRECTION ON CONVECTION VELOCITY = YES} (logical type, set {\ttfamily = NO} by default) modifies the depth-averaged convection velocity to account for the vertical gradients of velocities and concentration.

\subsection{Settling lag correction}
Sediment transport exhibits temporal lags with flow due to flow and sediment velocity difference and bed development~\cite{wu2007computational, Miles96}. A scaling factor that accounts for both the settling velocity and the lag time is therefore required for the saturation concentration profile to adjust to changes in the flow.
The keyword \telkey{SETTLING LAG} (logical type, set to \texttt{NON} by default) allows to compute the bed exchange factor beta based on Miles~\cite{Miles96}. This keyword must be used with the Nikuradse friction law, prescribed in the \telemac{2D} steering file as \telkey{LAW OF BOTTOM FRICTION = 5}.

